一阶语言的\emph{语义}由其\emph{结构}给出。
它由一个\emph{论域}构成，且论域中的元素被指派给每个常项符号。
函数符号由函数来解释，关系符号由论域上的关系来解释。
从变元集合到论域的函数称为\emph{变元指派}。
\emph{满足}关系将结构、变元指派和公式联系起来；
$\Sat{M}{!A}[s]$通过对~$!A$ 的结构进行归纳来定义。
$\Sat{M}{!A}[s]$仅依赖于对~$!A$ 中实际出现符号的解释，
特别地，若~$!A$ 不含自由变元，则它与~$s$ 无关。
因此，如果~$!A$ 是一个语句，
那么只要对任意（或所有）~$s$ 都有 $\Sat{M}{!A}[s]$，就有 $\Sat{M}{!A}$。

满足关系是所有语义概念的基础。
如果一个语句在每个结构中都被满足，则称它是\emph{有效的}，记作 $\Sat{{}}{!A}$。
如果所有满足语句集~$\Gamma$ 中每个语句的结构~$\Struct{M}$ 也都满足~$!A$，
则称语句~$!A$ 是语句集~$\Gamma$ 的\emph{语义后承}，记作 $\Gamma \Entails !A$。
如果存在某个结构使得~$\Gamma$ 中的每个语句都被满足，则称集合~$\Gamma$ 是\emph{可满足的}，
否则称其为\emph{不可满足的}。
这些概念是相互关联的，例如，$\Gamma \Entails !A$ 当且仅当 $\Gamma \cup \{\lnot
!A\}$ 不可满足。